\chapter{Theoretical Foundations}
\label{ch:foundations}

This chapter establishes the mathematical scaffolding on which the experimental
framework rests.  We proceed from first principles: a measure-theoretic
probability space with a filtration formalizing the no-lookahead constraint,
the Kelly Criterion and its asymptotic optimality guarantee, Bayesian conjugate
updating, and the Multi-Armed Bandit formalism including UCB and Thompson
Sampling strategies.  All results stated here are standard; we provide
self-contained statements and proof sketches to anchor the computational
experiments of Chapter~\ref{ch:experiments} to rigorous theory.

% ---------------------------------------------------------------
\section{Probability Space and Filtration}
% ---------------------------------------------------------------

\begin{definition}[Filtered Probability Space]
A \emph{filtered probability space} is a tuple $(\Omega, \Filt, \{\Filt_t\}_{t \geq 0}, \Prob)$
where $(\Omega, \Filt, \Prob)$ is a probability space and
$\{\Filt_t\}$ is a filtration, i.e.\ an increasing sequence of
$\sigma$-algebras $\Filt_0 \subseteq \Filt_1 \subseteq \cdots \subseteq \Filt$.
\end{definition}

In our setting the filtration is generated by observed returns.  Let
$r_1, r_2, \ldots$ be the sequence of per-period returns.  We set
\[
    \Filt_t \;=\; \sigma(r_1, r_2, \ldots, r_t), \qquad \Filt_0 = \{\emptyset, \Omega\}.
\]
An agent's betting fraction at time $t$ must be $\Filt_{t-1}$-measurable,
meaning it can depend only on information available \emph{before} $r_t$ is
revealed.  This is the formal statement of the no-lookahead constraint.

\begin{definition}[Admissible Strategy]
A sequence of betting fractions $\{f_t\}_{t \geq 1}$ is \emph{admissible} if
$f_t$ is $\Filt_{t-1}$-measurable for every $t$, and $f_t \in [0,1]$
almost surely.
\end{definition}

The wealth process under an admissible strategy satisfies the recursion
\begin{equation}
    W_{t} \;=\; W_{t-1}\,\bigl(1 + f_t\,r_t\bigr), \qquad W_0 = 1,
    \label{eq:wealth_recursion}
\end{equation}
which expanded telescopically gives
\[
    W_T \;=\; \prod_{t=1}^{T} \bigl(1 + f_t\,r_t\bigr).
\]
Taking logarithms converts this product to a sum, making log-wealth the
natural object of analysis:
\[
    \log W_T \;=\; \sum_{t=1}^{T} \log\!\bigl(1 + f_t\,r_t\bigr).
\]
\begin{remark}
All numerical computations in this thesis use $\log(1 + f_t r_t)$ via
\texttt{numpy.log1p} to avoid catastrophic cancellation when $f_t r_t$ is
small.  Ruin is defined as the absorbing event $1 + f_t r_t \leq 0$, at which
point $\log W_t = -\infty$.
\end{remark}

% ---------------------------------------------------------------
\section{The Kelly Criterion and Asymptotic Optimality}
% ---------------------------------------------------------------

\subsection{Binary Kelly Formula}

Consider a single bet with payoff odds $b > 0$: a wager of fraction $f$
returns $bf$ with probability $p$ and loses $f$ with probability $q = 1 - p$.
The expected log-growth per step is
\[
    G(f) \;=\; p\log(1+bf) + q\log(1-f).
\]

\begin{proposition}[Binary Kelly Fraction]
\label{prop:binary_kelly}
The unique maximizer of $G(f)$ over $f \in [0,1]$ is
\begin{equation}
    f^* \;=\; \frac{bp - q}{b} \;=\; p - \frac{q}{b},
    \label{eq:binary_kelly}
\end{equation}
provided $bp > q$ (positive expected value); otherwise $f^* = 0$.
\end{proposition}

\begin{proof}
$G$ is strictly concave on $[0,1)$ with derivative
$G'(f) = pb/(1+bf) - q/(1-f)$.
Setting $G'(f) = 0$ and solving for $f$ yields \eqref{eq:binary_kelly}.
Strict concavity of $G$ on $(0,1)$ follows since $G''(f) < 0$ for all
$f \in (0,1)$, guaranteeing uniqueness of the global maximum.
\end{proof}

\subsection{Continuous-Return Kelly Formula}

When returns are drawn from a continuous distribution, the Kelly fraction
maximizes $\E[\log(1 + fr)]$ over $f$.  For the Gaussian case
$r \sim \Normal(\mu, \sigma^2)$, a second-order Taylor expansion of
$\log(1+fr)$ around $f=0$ gives the approximation:
\begin{equation}
    f^*_{\text{Gaussian}} \;\approx\; \frac{\mu}{\sigma^2}.
    \label{eq:gaussian_kelly}
\end{equation}
This expression is exact when $|fr| \ll 1$ and provides a tractable closed
form for the oracle benchmarks used in Chapter~\ref{ch:experiments}.

\subsection{Strong Law of Large Numbers and Asymptotic Optimality}

The following theorem, due to Kelly~(1956) and rigorously formalized by
Breiman~(1961), is the central justification for maximizing log-utility.

\begin{theorem}[Kelly--Breiman Asymptotic Optimality]
\label{thm:kelly_slln}
Let $\{r_t\}_{t=1}^\infty$ be i.i.d.\ random variables with
$\E[\log(1+f^* r)] > -\infty$, and let
$f^* = \arg\max_{f \in [0,1]} \E[\log(1+fr)]$.
\begin{enumerate}
    \item \textbf{(Asymptotic Growth Rate)}
          \[
              \frac{\log W_T^{(f^*)}}{T} \;\xrightarrow{\text{a.s.}}\;
              \E\!\bigl[\log(1+f^*r)\bigr] \quad \text{as } T\to\infty.
          \]
    \item \textbf{(Dominance)} For any admissible strategy $\{f_t\}$ with
          $\liminf_{T} \frac{1}{T}\sum_t \log(1+f_t r_t) < \E[\log(1+f^* r)]$,
          \[
              \frac{W_T^{(f^*)}}{W_T^{(f)}} \;\xrightarrow{\text{a.s.}}\; +\infty.
          \]
\end{enumerate}
\end{theorem}

\begin{proof}[Proof sketch]
\textit{Part (1).}
By the Strong Law of Large Numbers applied to the i.i.d.\ sequence
$Z_t = \log(1 + f^* r_t)$:
\[
    \frac{1}{T}\sum_{t=1}^T Z_t \;\xrightarrow{\text{a.s.}}\; \E[Z_1]
    = \E\!\bigl[\log(1+f^*r)\bigr].
\]
Since $\frac{1}{T}\log W_T^{(f^*)} = \frac{1}{T}\sum_t Z_t$, part (1) follows.

\textit{Part (2).}
The difference in log-wealths satisfies
\[
    \frac{1}{T}\bigl(\log W_T^{(f^*)} - \log W_T^{(f)}\bigr)
    \xrightarrow{\text{a.s.}}
    \E[\log(1+f^*r)] - \liminf_T\frac{1}{T}\sum_t \log(1+f_t r_t) > 0,
\]
so $\log W_T^{(f^*)} - \log W_T^{(f)} \to +\infty$ a.s., implying the ratio
diverges.
\end{proof}

\begin{remark}[Finite-Horizon Cost]
Theorem~\ref{thm:kelly_slln} is asymptotic.  Over a finite horizon $T$,
Kelly betting can suffer severe drawdowns since the distribution of $W_T$
has a heavy right tail but also meaningful left-tail mass.  This motivates
the CPPI risk constraint introduced in Chapter~\ref{ch:methodology}.
\end{remark}

% ---------------------------------------------------------------
\section{Bayesian Updating and Posterior Beliefs}
% ---------------------------------------------------------------

\begin{definition}[Beta--Bernoulli Conjugate Pair]
Let outcomes $X_t \in \{0,1\}$ be i.i.d.\ Bernoulli$(p)$ with unknown $p$.
The \emph{Beta prior} $p \sim \text{Beta}(\alpha_0, \beta_0)$ is conjugate:
after observing $n$ wins and $m$ losses, the posterior is
\begin{equation}
    p \mid X_{1:n+m} \;\sim\; \text{Beta}(\alpha_0 + n,\; \beta_0 + m).
    \label{eq:beta_posterior}
\end{equation}
\end{definition}

The posterior mean and variance are
\[
    \hat p \;=\; \frac{\alpha_0 + n}{\alpha_0 + \beta_0 + n + m},
    \qquad
    \Var[p \mid \text{data}] = \frac{\hat p(1-\hat p)}{\alpha_0+\beta_0+n+m+1}.
\]
As $n + m \to \infty$, $\hat p \to p_{\text{true}}$ and the posterior
variance shrinks to zero, recovering the true parameter.

\begin{proposition}[Sticky Prior Effect]
\label{prop:sticky_prior}
Suppose an agent accumulates $T_0$ observations in a stationary regime with
true win rate $p_0$, so that $\alpha = \alpha_0 + n$, $\beta = \beta_0 + m$
with $n + m = T_0$ and $\hat p \approx p_0$.  If the true win rate then
switches to $p_1 \neq p_0$ at step $T_0 + 1$, the posterior mean requires
approximately
\[
    T_{\text{lag}} \;\approx\;
    \frac{\alpha_0 + \beta_0 + T_0}{|p_1 - p_0|}
\]
additional observations before it shifts to within $|p_1 - p_0|/2$ of $p_1$.
\end{proposition}

\begin{proof}
After $T_{\text{lag}}$ additional steps in the new regime the posterior mean is
\[
    \hat p(T_0 + T_{\text{lag}}) \;\approx\;
    \frac{(\alpha_0 + T_0 p_0) + T_{\text{lag}} p_1}{\alpha_0+\beta_0+T_0+T_{\text{lag}}}.
\]
Setting this equal to $(p_0 + p_1)/2$ and solving for $T_{\text{lag}}$ yields the
stated expression (to leading order in $T_0$).
\end{proof}

\begin{remark}
Proposition~\ref{prop:sticky_prior} makes the detection-lag problem precise.
An agent that has observed $T_0 = 200$ steps of bull-market data needs
$O(T_0) = O(200)$ additional observations before its posterior mean moves
meaningfully toward the bear-market rate.  The Volatility-Augmented HMM
(Chapter~\ref{ch:methodology}) bypasses this bottleneck by introducing a
second observational dimension --- rolling volatility --- that spikes
immediately upon a regime shift, triggering rapid reallocation of belief mass.
\end{remark}

% ---------------------------------------------------------------
\section{Multi-Armed Bandits}
% ---------------------------------------------------------------

\begin{definition}[Multi-Armed Bandit]
A $K$-armed bandit is a tuple $(\mathcal{A}, \{P_k\}_{k=1}^K)$ where
$\mathcal{A} = \{1,\ldots,K\}$ is the action set and $P_k$ is the unknown
reward distribution for arm $k$ with mean $\mu_k$.  At each step $t$ the
agent selects arm $A_t$ and observes reward $R_t \sim P_{A_t}$.
\end{definition}

Performance is measured by \emph{cumulative pseudo-regret}:
\[
    \mathcal{R}_T \;=\; T\mu^* - \sum_{t=1}^T \mu_{A_t},
    \qquad \mu^* = \max_k \mu_k.
\]

\subsection{Upper Confidence Bound (UCB1)}

\begin{theorem}[UCB1 Regret Bound, Auer et al.~2002]
\label{thm:ucb_regret}
The UCB1 algorithm selects $A_t = \arg\max_k
\left[\hat\mu_k + \sqrt{2\ln t\,/\,N_k(t)}\right]$
and achieves expected regret
\[
    \E[\mathcal{R}_T] \;\leq\; \sum_{k:\,\Delta_k > 0}
    \frac{8\ln T}{\Delta_k} + \left(1 + \frac{\pi^2}{3}\right)\sum_k \Delta_k,
\]
where $\Delta_k = \mu^* - \mu_k$.  This bound is $O(\log T)$ and
asymptotically matches the Lai--Robbins (1985) lower bound.
\end{theorem}

\subsection{Thompson Sampling}

Thompson Sampling (Thompson, 1933) selects arm $k$ at step $t$ with
probability equal to the posterior probability that arm $k$ is optimal:
\[
    A_t \sim \Prob(k = \arg\max_j \theta_j \mid \Filt_{t-1}),
\]
implemented by sampling $\tilde\theta_k \sim \text{Beta}(\alpha_k, \beta_k)$
and choosing $A_t = \arg\max_k \tilde\theta_k$.

\begin{theorem}[Thompson Sampling Regret, Agrawal and Goyal 2012]
\label{thm:ts_regret}
For Bernoulli bandits, Thompson Sampling achieves
$\E[\mathcal{R}_T] = O\!\left(\sqrt{KT\ln T}\right)$
and satisfies an asymptotically optimal instance-dependent bound matching the
Lai--Robbins lower bound up to constants.
\end{theorem}

In our framework, Thompson Sampling is combined with the Kelly formula:
after sampling $\tilde p_k \sim \text{Beta}(\alpha_k,\beta_k)$, the bet size
for arm $k$ is set to $f^*(\tilde p_k)$ from \eqref{eq:binary_kelly}.
This preserves the Bayesian exploration property while sizing bets to the
sampled belief, providing an implicit risk-aversion mechanism: high posterior
variance produces a wide spread of $\tilde p_k$ draws, leading to variable
and on average smaller Kelly fractions than a posterior-mean plug-in.

\subsection{EXP3 for Adversarial Bandits}

EXP3 (Auer et al., 2002b) is designed for the adversarial setting where the
reward sequence can be adversarially chosen.

\begin{theorem}[EXP3 Regret Bound]
\label{thm:exp3_regret}
EXP3 with parameter $\gamma \in (0,1]$ achieves expected regret against any
adaptive adversary:
\[
    \E[\mathcal{R}_T] \;\leq\; (e-1)\gamma G^* + \frac{K\ln K}{\gamma},
\]
where $G^* = \max_k \sum_t r_{k,t}$.  Setting
$\gamma = \min\!\left(1, \sqrt{K\ln K\,/\,((e-1)T)}\right)$
yields $\E[\mathcal{R}_T] = O(\sqrt{TK\ln K})$.
\end{theorem}

\begin{remark}[Implementation Note --- Fractional EXP3]
\label{rem:exp3_variant}
Canonical EXP3 selects a single arm each round.  Our implementation uses the
EXP3 probability weights directly as portfolio fractions, scaled by a base
leverage, yielding a continuous-fraction variant we call \emph{EXP3 (fractional)}.
This variant does not satisfy Theorem~\ref{thm:exp3_regret} exactly; the bound
provides a motivation for its robustness properties rather than a strict
guarantee.  All results in Chapter~\ref{ch:experiments} label this agent
accordingly.
\end{remark}

% ---------------------------------------------------------------
\section{Regime-Switching Environments and the HMM}
% ---------------------------------------------------------------

\begin{definition}[Hidden Markov Model]
A discrete-time HMM consists of:
(i) a finite hidden state space $\mathcal{S} = \{s_1,\ldots,s_M\}$;
(ii) an initial distribution $\pi$;
(iii) a transition matrix $A$ with $A_{ij} = \Prob(S_t = s_j \mid S_{t-1} = s_i)$;
(iv) emission distributions $b_j(\cdot) = p(\cdot \mid S_t = s_j)$.
The observations $\{y_t\}$ are conditionally independent given the states.
\end{definition}

The optimal filter $\Prob(S_t \mid y_{1:t})$ is computed by the
\emph{forward algorithm}.  In log-space for numerical stability:
\begin{equation}
    \ln\alpha_t(j) \;=\; \ln b_j(y_t)
    + \operatorname{logsumexp}_i\!\bigl[\ln\alpha_{t-1}(i) + \ln A_{ij}\bigr],
    \label{eq:forward_logspace}
\end{equation}
normalized via $\ln\alpha_t \leftarrow \ln\alpha_t - \operatorname{logsumexp}(\ln\alpha_t)$.

\begin{remark}[Specified vs.\ Learned Emissions]
\label{rem:hmm_specified}
In Chapter~\ref{ch:methodology} the emission parameters
$(\mu_j, \sigma_j, k_j, \theta_j)$ are specified a priori based on domain
knowledge of bull and bear return distributions, rather than estimated via
the Baum--Welch EM algorithm.  With only $T \approx 250$ observations,
EM can overfit the emission model; the a priori specification yields a more
robust filter.  This framework is therefore a \emph{Bayesian filter with a
specified generative model} rather than a fully learned HMM.
\end{remark}
